% ID: SEL-D0001-OA-E04
% ESTADO: cotejada. Fuente y transcripción revisadas de forma independiente; salida compilada inspeccionada.
% Fuente: Selectividad/sel_2019_matematicas_aplicadas/Titular Junio/Exámenes Mat Apli CCSS A y B.pdf
% SHA-256: fa0cbbddfdf34b1f9fcf8d1090c2d0c3f13e5353e06da22ec354d3345e313584
% Página: 1; opción A; ejercicio 4.
% Encabezado original: EJERCICIO 4. Solo enunciado; no contiene solución.

Se desea estimar la proporción de individuos que piensan votar a un cierto
partido político en una determinada ciudad. Para ello se toma una muestra
aleatoria de 300 individuos de la ciudad, resultando que 135 de ellos piensan
votar a ese partido.

\begin{enumerate}[label=\alph*)]
  \item \textbf{(1.5 puntos)} Calcule un intervalo de confianza al 97\,\% para la
  proporción de individuos que piensan votar a ese partido en dicha ciudad.
  \item \textbf{(1 punto)} Suponiendo que se mantiene la misma proporción
  muestral y el mismo nivel de confianza del apartado anterior, determine
  el tamaño mínimo de la muestra para estimar la proporción con un error
  inferior al 2\,\%.
\end{enumerate}
